%==============================================================================
%  4. Experimental Results and Discussion
%==============================================================================

\section{Experimental Results and Discussion}
\label{sec:experiments}

\subsection{Experimental Setup}
\label{subsec:setup}

\textbf{Scenarios.} Two scenarios are constructed over a DEM of the Chengguan
district of Lhasa ($91.02^{\circ}$--$91.28^{\circ}$E,
$29.52^{\circ}$--$29.73^{\circ}$N), with two radar threat fields of radius
\SI{5000}{m} and \SI{6000}{m}. The \emph{balanced} scenario has
$\Np = \Mt = 10$, so every UAV can be assigned a distinct target; the
\emph{scarce-resource} scenario has $\Np = 15$ and $\Mt = 30$, which tightens the
assignment constraints considerably.

\textbf{Baseline and variants.} Two configurations form the main ablation:
a vanilla \dmde{} baseline (A0) in which \CR{}, \Ffactor{} and \GMR{} are
coupled by formulas~3-9, 3-11 and 3-12, and the full framework (A1,
``+Search Controller'') in which they are decoupled. Three further control
conditions isolate the action space: decoupled (A1), coupled (\CR{} only) and
CR-locked (\Ffactor{} and \GMR{} only).

\textbf{Parameters.} Population size $\NP = 50$; maximum generations
$\Gmax = 1000$; scaling factor $\Ffactor = 0.5$;
\todo{verify the baseline values of \Ffactor{} and \CR{}};
30 independent runs per configuration with seeds $42$--$71$. The LLM is
Qwen3-14B served by SiliconFlow with thinking disabled, temperature $0.7$, and
\searchctrl{} is queried every $I = 50$ generations. In the CR-locked condition
\CR{} is pinned to $0.3$, the best value found by an offline sweep.

\textbf{Metrics.} Solution quality (best, mean $\pm$ standard deviation and
median of the final cost), convergence speed (generations needed to reach $90\%$,
$95\%$ and $99\%$ of the total improvement), and wall-clock cost decomposed into
the DMDE loop and LLM calls.

% 待补：硬件环境（CPU 型号/内存）、软件版本与代码可获取性说明
\todo{report the hardware (CPU, memory), software versions and how to obtain
      the code.}

\subsection{Comparison with Existing Methods}
\label{subsec:comparison}

% 这一节与 4.3 的边界：4.2 与外部方法/基线整体对比，4.3 做内部消融

% 待补：填入与基线方法的整体对比结果与讨论。若暂时只做内部消融，
% 可把本节与 4.3 合并，并在 README 中同步更新章节编号
\todo{insert the comparison against external baselines and its discussion;
      if only the internal ablation is available, merge this subsection
      into 4.3 and update the section numbering in README.md.}

\subsection{Ablation Study}
\label{subsec:ablation}

\tabref{tab:ablation_results} reports the solution quality of the main
configurations on both scenarios, and \figref{fig:convergence} shows the
corresponding convergence curves.

% --- 表格：从 experiments/ 同步而来（见 tables/README.md）---
% \input{tables/ablation_solution_quality}

\begin{table}[htbp]
  \centering
  \caption{Solution quality over 30 runs. Best values are shown in bold where a
           tie occurs.}
  \label{tab:ablation_results}
  \begin{tabular}{llcccc}
    \toprule
    Scenario & Configuration & Best & Mean $\pm$ Std & Median \\
    \midrule
    balanced & Vanilla \dmde{}      & \dots & \dots & \dots \\
             & + Search Controller  & \dots & \dots & \dots \\
    \midrule
    scarce   & Vanilla \dmde{}      & \dots & \dots & \dots \\
             & + Search Controller  & \dots & \dots & \dots \\
    \bottomrule
  \end{tabular}
\end{table}

\begin{figure}[htbp]
  \centering
  % \includegraphics[width=0.8\linewidth]{convergence_S1.pdf}
  \fbox{\parbox[c][3.5cm][c]{0.8\linewidth}{\centering TODO: convergence curves}}
  \caption{Mean convergence curves on the balanced scenario (shaded band:
           $\pm$ one standard deviation over 30 runs).}
  \label{fig:convergence}
\end{figure}

% 待补：讨论要点——三个动作空间条件的差异；CR 决策的自相关性（实测
% corr(CR_t, CR_{t-1}) = -0.967、翻转率 0.98）说明模型在拟合噪声；
% 相邻档位效应 -0.248pp vs 运行间噪声 sd 3.513pp，信噪比约 1/14；
% 以及显著性检验结论（离线固定 CR=0.3 在 9/9 组对照中胜出）
\todo{discuss: the difference between the three action-space conditions; the
      autocorrelation of the \CR{} decisions (measured
      $\mathrm{corr}(\CR_t, \CR_{t-1}) = -0.967$ with a flip rate of $0.98$),
      which indicates that the model is fitting noise; the effect between
      adjacent \CR{} levels ($-0.248$~pp) against the between-run noise
      ($3.513$~pp), i.e. a signal-to-noise ratio of about $1/14$; and the
      significance-test outcome, in which an offline constant $\CR = 0.3$
      wins 9 out of 9 comparisons.}

\subsection{Scalability Analysis}
\label{subsec:scalability}

% 待补：随 N/M 增长时的解质量与 LLM 调用开销趋势，说明可扩展性与预算增长关系
\todo{report how solution quality and LLM query overhead scale with $\Np$ and
      $\Mt$, and whether the query budget grows linearly or sub-linearly.}
